% Generated draft: requires mathematical and source review.
\documentclass[11pt]{article}
\usepackage[margin=25mm]{geometry}
\usepackage{fontspec}
\setmainfont{Latin Modern Roman}
\usepackage{amsmath,amssymb,mathtools,unicode-math}
\setmathfont{Latin Modern Math}
\usepackage{enumitem,needspace,xurl,hyperref}
\setlength{\parindent}{0pt}
\setlength{\parskip}{5pt}
\setlength{\emergencystretch}{3em}
\newcommand{\N}{\mathbb{N}}
\newcommand{\Z}{\mathbb{Z}}
\newcommand{\Q}{\mathbb{Q}}
\newcommand{\R}{\mathbb{R}}
\newcommand{\C}{\mathbb{C}}
\newcommand{\F}{\mathbb{F}}
\newcommand{\A}{\mathbb{A}}
\newcommand{\PP}{\mathbb{P}}
\newcommand{\E}{\mathbb{E}}
\newcommand{\eps}{\varepsilon}
\newcommand{\dd}{\,\mathrm{d}}
\begin{document}
\section*{UnsolvedMath ID: 30002980}
\textbf{Source problem code:} catalogue entry
\subsection{Definitions and mathematical statement}
The supplied record asks about Brill--Noether intersections associated with a \(K3\) surface, but does not define the maps \(\psi_{X,n,\delta}\) and \(\psi_{p,n,\delta}\), their domains, or their parameter ranges. The following therefore preserves the question with its unresolved notation explicitly qualified.

For a map \(f\colon S\to T\), write \(\operatorname{Im}(f)=\{f(s):s\in S\}\). If both source maps take values in an ambient moduli space denoted by \(M_g\), and \(M^r_{g,d}\subseteq M_g\) is the intended Brill--Noether locus, put
\[
I_X=\operatorname{Im}(\psi_{X,n,\delta})\cap M^r_{g,d},
\qquad
I_p=\operatorname{Im}(\psi_{p,n,\delta})\cap M^r_{g,d}.
\]
These expressions are conditional on specifying the maps and a common ambient moduli space; neither specification is contained in the supplied record.

Under a usual convention, \(M_g\) parametrizes smooth projective connected curves of genus \(g\), and the underlying set of a Brill--Noether locus is
\[
M^r_{g,d}
 =\bigl\{[C]\in M_g:\text{there exists }L\in\operatorname{Pic}^{d}(C)
 \text{ with }h^0(C,L)\ge r+1\bigr\},
\]
where \(\operatorname{Pic}^{d}(C)\) denotes the set of isomorphism classes of degree-\(d\) line bundles and \(h^0(C,L)=\dim H^0(C,L)\). This is an explanatory convention, not a verified definition from the source report. It does not impose base-point freeness or require a complete linear series of dimension exactly \(r\); whether the intended locus imposes such conditions is unresolved. The ground field and the use of a coarse moduli space, stack, or compactification are also unspecified.

Likewise, \(X^{[d]}\) conventionally denotes the Hilbert scheme of zero-dimensional subschemes of \(X\) of length \(d\). This interpretation requires an appropriate surface \(X\) and a nonnegative integer \(d\); the record does not provide the precise hypotheses on \(X\).

\textbf{Source question.} With the notation and hypotheses intended in the cited report, describe \(I_X\) and \(I_p\), and determine what consequences the first intersection \(I_X\) has for the Mori theory of \(X^{[d]}\).

The record does not specify the meanings of \(p,n,\delta\), their relation to \(g\), the polarization or generality assumptions on \(X\), admissible parameter ranges, the intersection structure intended, or a particular Mori-theoretic conclusion. Consequently, it does not determine a fully quantified mathematical conjecture. These missing qualifications must be recovered before this draft can serve as a precise standalone formulation.
\subsection{Short English statement}
The question asks how the images of two unspecified moduli maps meet a Brill--Noether locus, and what the first intersection implies for the Mori theory of the Hilbert scheme \(X^{[d]}\). Essential definitions and hypotheses are missing from the supplied record, so the formulation remains conditional and requires human mathematical review.
\subsection{Sources}
\begin{itemize}
\item Local registry: \texttt{data/unsolvedmath/problems.json}, with \texttt{catalog/problems.json} as the additional registry specified by the project; permanent UnsolvedMath ID \(30002980\), source problem code \texttt{catalogue entry}, title ``Brill--Noether Intersections of K3 Moduli Images.'' This draft uses the supplied complete record; the registry files were not independently inspected.
\item Supplied report attribution: \emph{Mini-Workshop: Singular Curves on \(K3\) Surfaces and Hyperkähler Manifolds} (2016), \emph{Oberwolfach Reports}, pp.~2939--2967, DOI \texttt{10.4171/owr/2015/51}. \url{https://doi.org/10.4171/owr/2015/51}. The registry gives proposed year 2015 and supplies no individual proposer.
\item Supplied literature reference: Ciliberto--Knutsen, \emph{k-Gonal Loci on K3 Surfaces}. \url{https://doi.org/10.1016/j.jsc.2013.11.007}.
\item Supplied literature reference: Knutsen--Torelli, \emph{Prescribed-Ramification Brill--Noether Loci}. \url{https://arxiv.org/2411.14562}.
\item Supplied literature reference: \emph{Singular Curves on a K3 Surface and Linear Series}. \url{https://citeseerx.ist.psu.edu/document?doi=204c4253d30a59a7601c8db2a18d4e49dc04c964&repid=rep1&type=pdf}. No authors are supplied for this entry.
\item Literature assessment dated \texttt{2026-08-21}: the supplied curation labels the problem ``partially solved'' and states that Brill--Noether loci in Severi varieties and their relation to rational curves in Hilbert schemes are substantially developed, while the requested intersections and resulting Mori-theoretic consequences are not classified in general. This is supplied dated curation, not an independently verified assessment. The background expressly calls for independent verification of time-sensitive status and cited claims.
\item The registry labels the statement \texttt{corrected\_verified} and reports that it was checked against the official Oberwolfach report. This is a supplied verification claim; neither that report nor the literature references were visited for this draft.
\end{itemize}
\end{document}
